\section{Mathematical Model Hierarchy}
\label{sec:model-hierarchy}

The central organizing device of the review is the model hierarchy. The
hierarchy is not a ladder of accuracy. It is a way to ask what a formulation
retains, what it closes, what it can observe natively, and what it must map
before a comparison is meaningful. The major review sources contribute
different parts of that frame. Quarteroni and Formaggia place cardiovascular
simulation inside a continuum and multiscale modeling program; Formaggia,
Lamponi, and Quarteroni identify the variables and closures retained by
distributed 1D artery models; Shi, Lawford, and Hose organize the 0D/1D network
literature around pressure--flow states and terminal coupling; and
Koppl--Helmig recast the same landscape as dimension reduction from detailed
flow fields to lower-dimensional states
\cite{QuarteroniFormaggia2004CardiovascularModeling,FormaggiaEtAl2003OneDimensionalBloodFlow,ShiEtAl2011BloodFlowModelReview,KopplHelmig2023DimensionReducedModeling}.
Read together, these sources motivate a hierarchy organized by retained state
and observation, not by a simple claim that higher dimension is always the
appropriate model.
Here \emph{retained state} means the variables advanced or constrained by the
mathematical formulation itself. A \emph{native observable} is computed
directly from that state before any reconstruction, diagnostic closure, or
cross-tier output map is applied.

The dimensional labels below therefore describe retained state and observation
contracts, not merely mesh geometry. A 0D model stores lumped pressure, volume,
or flow variables and is often used as a terminal or network closure for a
distributed or resolved domain. A 1D vessel model advances axial balance laws
for section area, flow, mean velocity, and wall-law pressure along a centerline
\cite{FormaggiaEtAl2003OneDimensionalBloodFlow,ShiEtAl2011BloodFlowModelReview,KopplHelmig2023DimensionReducedModeling}.
An axisymmetric or 2D model still solves spatial PDEs, but it retains variation
only in selected coordinates or imposes symmetry, suppressing the remaining
direction. Resolved 3D CFD advances velocity and pressure fields on the lumen,
while resolved FSI adds structural wall degrees of freedom and interface
traction or kinematic coupling
\cite{QuarteroniFormaggia2004CardiovascularModeling,FormaggiaEtAl2007FSICoupling}.
Thus 0D, 1D, and 2D tiers are not simply alternative ``3D CFD methods''; they
are reduced or coupled formulations with different native states and observation
maps.

\subsection{Resolved 3D CFD and FSI Models}
\label{subsec:resolved-reduced-model-tiers}

The resolved continuum tier starts from conservation of mass and momentum on a
three-dimensional lumen. In a fixed-wall CFD model the wall is a prescribed
geometric boundary, while in resolved FSI the wall has structural degrees of
freedom and transfers kinematic and traction information through a moving
interface
\cite{QuarteroniFormaggia2004CardiovascularModeling,FormaggiaEtAl2007FSICoupling}.
These models can expose local velocity gradients, wall traction, pressure
fields, recirculation, separation, and geometry-specific flow structures. A WSS
quantity, for example, is not just a scalar label: it requires a stress tensor,
wall normal, tangent projection convention, and surface output map
\cite{JohnEtAl2017WallShearStressVectorForm}.

The resolved tier also contains distinct uses of ``three-dimensional'' data.
Peskin's numerical analysis of blood flow in the heart is an early example in
which geometry, moving boundaries, and numerical representation are already
part of the mathematical problem, not merely input data
\cite{Peskin1977NumericalAnalysisBloodFlowHeart}. Quarteroni--Formaggia-style
cardiovascular modeling then broadens that perspective across continuum flow
and system coupling, while 3D--1D FSI coupling work makes the interface and
wall state explicit
\cite{QuarteroniFormaggia2004CardiovascularModeling,FormaggiaEtAl2007FSICoupling}.
By contrast, a fixed-wall resolved stenosis computation used to estimate a
pressure or velocity diagnostic may retain the same local velocity and pressure
fields while closing wall motion and downstream circulation differently. The
source-to-source contrast is therefore not simply 3D versus reduced. It is
which local fields are retained, which wall and boundary variables are closed,
and which observable is extracted from the resulting field.

\subsection{Axisymmetric and Two-Dimensional Reductions}
\label{subsec:axisymmetric-two-dimensional-models}

Axisymmetric and two-dimensional models occupy an intermediate role. They
retain more geometry and local velocity information than a one-dimensional
area--flow model, but they remove at least one spatial direction or impose
symmetry. Classical idealized stenosis studies use this controlled setting to
separate obstruction, pressure loss, and flow-structure effects from
patient-specific geometry and boundary uncertainty
\cite{YoungTsai1973SteadyStenosis,YoungTsai1973UnsteadyStenosis}. Their
contribution is therefore not only that they occupy a lower-dimensional
category; it is that they show how idealization can isolate a stenosis
mechanism.

The same reduction also imposes a limit. Axisymmetric and two-dimensional
models remain useful for sectional diagnostics and controlled comparisons only
when their boundary conditions, rheology, wall treatment, and observable
definitions are declared. Suppressed secondary flow, asymmetric geometry, and
wall mechanics cannot later be recovered as native outputs unless an additional
closure or observation map supplies them.

\subsection{Distributed One-Dimensional Area--Flow Models}
\label{subsec:distributed-one-dimensional-models}

Distributed 1D models reduce each vessel segment to axial variables, most often
cross-sectional area, flow, mean velocity, and pressure. Formaggia, Lamponi,
and Quarteroni make this state reduction concrete through area--flow equations,
wall laws, and profile assumptions, while Koppl--Helmig emphasize how such
models should be read as dimension-reduced representations with explicit
closure and coupling choices
\cite{FormaggiaEtAl2003OneDimensionalBloodFlow,KopplHelmig2023DimensionReducedModeling}.
This tier preserves pulse propagation and vessel-to-vessel transport at far
lower cost than resolved CFD/FSI, but profile, wall, friction, rheology,
stenosis, and boundary effects enter as closures rather than resolved fields.
It is therefore appropriate for single-vessel stenosis calculations when the
target observables are area, flow, mean velocity, and reduced pressure, and
when unavailable transverse flow, wall traction, and local stress fields are
not required as native outputs.

Stenosis-specific reductions also divide by observable. Pressure-loss formulas
and pressure-drop reduced-order models, such as the Lyras--Lee pressure-drop
ROM, target obstruction through a flow--loss relation
\cite{LyrasLee2021PressureDropROM}. Extended 1D stenosis formulations instead
place narrowed reference geometry and wall response inside an area--flow system
\cite{CanicEtAl2024Extended1DStenosis}. The Canic et al. arXiv preprint
reconstructs radial velocity only in post-processing, so that radial field is
not a native 1D state. Network-scale coronary models add tree
structure and terminal assumptions, so the Duanmu coronary-tree contribution is
not just another single-vessel stenosis model; it shows how reduced vessel
states become part of a larger pressure--flow network
\cite{DuanmuEtAl2019CoronaryArterialTree}. These sources should therefore be
compared by target observable and closure contract, not collected under a
single ``1D'' label.

\subsection{Zero-Dimensional Networks and Multiscale Coupling}
\label{subsec:network-multiscale-learned-models}

Zero-dimensional and network models replace distributed fields with lumped
states and algebraic or ordinary-differential pressure--flow relations. The
Shi--Lawford--Hose review is useful here because it treats 0D and 1D models as
circulation models with explicit resistance, compliance, junction, and terminal
assumptions, rather than as merely cheaper versions of resolved CFD
\cite{ShiEtAl2011BloodFlowModelReview}. Their state variables are not local
wall or velocity fields. In a single-vessel calculation, these models most
often enter through inlet or outlet closure rather than as the main spatial
solver.

Geometrical multiscale models connect tiers, for example 3D--1D, 3D--0D, or
1D--0D, through interface conditions that transfer pressure, flow, traction,
impedance, or characteristic information
\cite{QuarteroniVeneziani2003GeometricalMultiscale,FormaggiaEtAl2007FSICoupling}.
Quarteroni and Veneziani foreground the ODE--PDE coupling problem, while
Formaggia, Moura, and Nobile expose the additional stability and interface
issues that arise when the coupled tiers include FSI. The mathematical issue is
therefore not only that two solvers are coupled; it is that the interface
variables must have compatible meanings across dimensions. Coronary FFR
comparisons between 1D and 3D models show the same issue on the output side:
the reduced and resolved states must be mapped to a common pressure-ratio
observable before agreement or disagreement has a mathematical meaning
\cite{Blanco2018FFR1D3D}.

\subsection{Learned Surrogates, PINNs, and Operator Maps}
\label{subsec:learned-surrogates-operator-maps}

Surrogate and scientific-ML models add a second axis to the hierarchy. The
general PINN framework constrains an approximate solution or inverse map with
residual, boundary, data, and regularization terms, while hemodynamics-specific
reviews ask how that framework translates to vascular flow settings
\cite{RaissiEtAl2019PhysicsInformedNeuralNetworks,YuEtAl2026PINNHemodynamicsReview}.
DeepONet and Fourier neural operators shift the claim from one computed
solution to a learned map between functions or parameterized fields
\cite{LuEtAl2021DeepONet,LiEtAl2021FNO}. Hemodynamics-specific neural
differential and deep-operator examples then target reduced blood-flow or
stenosed-vessel maps directly. Csala and coauthors report conditional error
reductions in an idealized stenosis/waveform setting, while Velikorodny and
coauthors derive FFR after learning pressure and flow maps in a simplified
stenosis setting \cite{CsalaEtAl2025PCNDE,VelikorodnyEtAl2025DeepOperatorStenosed}.
These learned categories do not remove the need to state the underlying model
tier: every learned-model claim should identify its training target, input
functions or parameters, output observable, geometry family, and evaluation
domain.

\subsection{Native States and Cross-Tier Observation Maps}
\label{subsec:native-states-observation-maps}

Tables~\ref{tab:model-hierarchy-comparison} and
\ref{tab:model-hierarchy-observation-evidence} summarize the hierarchy in two
steps. The first table records retained state, reduction, closure, and data
requirements. The second records observables, evidence role, and interpretation
limits. The split keeps the reading frame organized around state and
observation rather than around individual papers.

\begin{table}[!htb]
    \centering
    \scriptsize
    \caption{State, reduction, closure, and data contract for cardiovascular
    model tiers.}
    \label{tab:model-hierarchy-comparison}
    \begin{tabularx}{\textwidth}{@{}
        >{\raggedright\arraybackslash}p{0.15\textwidth}
        >{\raggedright\arraybackslash}p{0.22\textwidth}
        >{\raggedright\arraybackslash}p{0.16\textwidth}
        >{\raggedright\arraybackslash}p{0.21\textwidth}
        >{\raggedright\arraybackslash}X@{}}
        \toprule
        \textbf{Tier}
            & \textbf{Native state and form}
            & \textbf{Reduction}
            & \textbf{Closure burden}
            & \textbf{Data or coupling requirement} \\
        \midrule
        Fixed-wall 3D CFD
            & $\vu,p$ on a prescribed lumen; Navier--Stokes or Stokes PDE.
            & None in space.
            & Rheology, wall boundary, numerics.
            & Geometry and inlet/outlet data. \\
        Resolved FSI
            & Fluid fields plus wall state; coupled fluid/structure PDEs.
            & Structural idealization only.
            & Wall material, interface, mesh/ALE.
            & Wall data and coupling histories. \\
        Axisymmetric / 2D
            & Symmetry-constrained velocity and pressure; reduced spatial PDE.
            & Suppressed azimuthal/asymmetric modes.
            & Symmetry, wall, rheology, boundaries.
            & Radius/profile data. \\
        Distributed 1D
            & $A,Q,\bar u$, pressure/wall state; axial balance law.
            & Section averaging.
            & Profile, wall, friction, geometry, rheology.
            & Inlet/outlet, junctions, terminal beds. \\
        Lumped 0D
            & Pressures, volumes, flows; ODE/algebraic network.
            & Spatial distribution removed.
            & Resistance, compliance, inertance, valves.
            & Network topology and terminal data. \\
        Geometrical multiscale
            & Tier-specific states plus interfaces; coupled PDE/ODE systems.
            & Local/system split.
            & Interface compatibility and stability.
            & Pressure, flow, traction, impedance. \\
        Learned surrogate
            & Learned solution, parameter, or operator map; trained objective.
            & Inherits source model/data tier.
            & Architecture, loss, training distribution.
            & Training inputs, geometry family, labels. \\
        \bottomrule
    \end{tabularx}
\end{table}

\begin{table}[!htb]
    \centering
    \scriptsize
    \caption{Observable, evidence, and interpretation limits for cardiovascular
    model tiers.}
    \label{tab:model-hierarchy-observation-evidence}
    \begin{tabularx}{\textwidth}{@{}
        >{\raggedright\arraybackslash}p{0.16\textwidth}
        >{\raggedright\arraybackslash}p{0.22\textwidth}
        >{\raggedright\arraybackslash}p{0.22\textwidth}
        >{\raggedright\arraybackslash}X@{}}
        \toprule
        \textbf{Tier}
            & \textbf{Native observables}
            & \textbf{Reconstructed or diagnostic outputs}
            & \textbf{Evidence role and limit} \\
        \midrule
        Fixed-wall 3D CFD
            & Velocity and pressure fields.
            & WSS indices and section means.
            & Supports local-field verification; wall motion and circulation
            are closed. \\
        Resolved FSI
            & Fluid fields, wall displacement, and interface traction.
            & Scalar WSS or pressure diagnostics.
            & Supports coupled-interface analysis; many data and material
            inputs must be identified. \\
        Axisymmetric / 2D
            & Symmetry-constrained sectional fields.
            & 3D asymmetry proxies.
            & Supports mechanism isolation; suppressed modes are not native. \\
        Distributed 1D
            & Area, flow, mean velocity, and wall-law pressure.
            & WSS proxies and radial-profile reconstructions.
            & Supports pulse or stenosis comparison; local traction and
            transverse flow are closed. \\
        Lumped 0D
            & Global pressure--flow states.
            & Local fields only by added model.
            & Supports system interaction; stenosis geometry and WSS are not
            native. \\
        Geometrical multiscale
            & Tier-native and interface values.
            & Cross-tier diagnostics.
            & Supports coupled-domain consistency; interface maps can dominate
            interpretation. \\
        Learned surrogate
            & Declared learned target.
            & Derived diagnostics such as FFR.
            & Supports in-distribution generalization tests; credibility is
            bounded by the data-generating contract. \\
        \bottomrule
    \end{tabularx}
\end{table}

Cross-tier comparisons are therefore observed-output comparisons:
\begin{flalign*}
    \quad
    Y_{1D}=\mathcal O_{1D}(U_{1D}),
    \qquad
    Y_{3D}=\mathcal O_{3D}(\vu,p,\OmgB),
    \qquad
    \delta Y=Y_{1D}-Y_{3D}. &&
\end{flalign*}
Unless the comparison design separates them, $\delta Y$ combines model-form,
closure, parameter, numerical, data-matching, and observation-operator effects.
Blanco's 1D--3D FFR study is read here as an FFR-specific matched-contract
comparison with pressure boundary matching and dedicated stenosis/bifurcation
losses, not as evidence of general model-tier equivalence
\cite{Blanco2018FFR1D3D}.

\begin{stenosisillustration}
\label{ex:stenosisillustration-model-hierarchy}
The case study later in the report sits in the distributed 1D branch. Its
mathematical state is a single-vessel area--flow system with an analytic
stenosis reference geometry and reduced wall/friction/profile closures. The
case therefore belongs beside other 1D stenosis models, including extended
variable-radius formulations \cite{CanicEtAl2024Extended1DStenosis}, rather
than beside resolved CFD or FSI as a native wall-traction or wall-displacement
solver. When it is compared with resolved data, the comparison passes through
an observation map from 1D area--flow variables to the chosen velocity, flow, or
pressure quantity.
\end{stenosisillustration}
